\subsection{Parejas Incompatibles en Secuencias de Paréntesis (Bloques de Nivel Superior)}

\graybold{Preguntas guía:}

\begin{itemize}
\item ¿Trabajo con una secuencia balanceada de paréntesis y necesito razonar sobre qué emparejamientos específicos son válidos?
\item ¿La pregunta involucra qué pasa con el RESTO de la secuencia al fijar/quitar un par específico?
\item ¿Puedo identificar los "bloques de nivel superior" (tramos donde el balance nunca vuelve a cero hasta el final del bloque)?
\end{itemize}

\graybold{¿Para qué sirve?} Analizar, para una secuencia balanceada de paréntesis, la validez de emparejamientos NO necesariamente anidados de la forma estándar, descomponiendo la secuencia en bloques de nivel superior (tramos maximales donde el balance se mantiene estrictamente positivo hasta el cierre). La clave es que un '(' y un ')' se pueden emparejar sin romper la validez del resto si y solo si están en el MISMO bloque.

\graybold{Complejidad:} \bigO{n}.

\graybold{Fundamento teórico:} Sea $\text{bal}(p)$ el balance (número de '(' menos número de ')') tras procesar los primeros $p$ caracteres. Un bloque de nivel superior termina exactamente donde $\text{bal}(p) = 0$. Al emparejar una posición $i$ ('(') con una posición $j$ (')'), $i<j$, el balance de las posiciones ENTRE $i$ y $j-1$ baja en 1 respecto al original (porque ya se "usó" el '(' pero no el ')'); esto rompe la validez del resto si y solo si el balance ORIGINAL en algún punto de ese rango ya era 0 (pues bajaría a -1). Como los únicos puntos con balance 0 son los cierres de bloque, el par es INCOMPATIBLE si y solo si $i$ y $j$ pertenecen a bloques distintos. El total de parejas incompatibles se calcula entonces como (total de parejas $(i,j)$ válidas en toda la secuencia) menos (parejas válidas contando solo dentro de cada bloque por separado). Ambas cantidades se obtienen en un único recorrido llevando un contador de '(' abiertos, global y reiniciado por bloque.

\graybold{Ejercicio aplicado}

Dada una secuencia de N bailarines Yang '(' y Yin ')' (garantizada balanceada), donde cada Yang debe emparejarse con un Yin posterior, contar cuántas parejas $(i,j)$ son "incompatibles": aquellas que, de formarse, impiden que el resto de bailarines encuentre pareja.

\graybold{I/O:} N ≤ 10\textsuperscript{6} → un solo recorrido lineal con lectura total (patrón 2); no hace falta tokenizar por espacios ya que la secuencia es un único string contiguo. Verificado: \textasciitilde{}0.1s en el peor caso.

\graybold{Código solución}

\begin{lstlisting}[language=Python]
import sys

def solve():
    data = sys.stdin.buffer.read().split()
    n = int(data[0])
    s = data[1].decode()

    total_pairs = 0     # total de parejas (i,j) con i<j, s[i]='(', s[j]=')'
    within_block = 0    # de esas, cuantas quedan DENTRO del mismo bloque de nivel superior
    open_global = 0
    open_block = 0
    bal = 0

    for ch in s:
        if ch == '(':
            open_global += 1
            open_block += 1
            bal += 1
        else:
            total_pairs += open_global
            within_block += open_block
            bal -= 1
            if bal == 0:          # cierre de un bloque de nivel superior y reinicia el contador de bloque
                open_block = 0

    print(total_pairs - within_block)

solve()
\end{lstlisting}